\section{Explicit check of the 11d to 10d Killing spinor reduction}
\label{sec:IIA killing spinor reduction}
We verify explicitly that a covariantly-constant spinor of 11d supergravity reduces, on the standard Kaluza-Klein ansatz, to a Killing spinor of the 10d type IIA theory in the gravity, KK-vector and dilaton sector, reproducing the gravitino variation \eqref{IIA Killing spinor equation}.
We label the eleventh coordinate $x^{10}$ and its flat (frame) index likewise by $10$; $\Gamma^{10}$ is then the 10d chirality operator usually written $\Gamma^{11}$.

The reduction ansatz is
\begin{equation}
G_{MN}
=e^{-2\phi/3}
\begin{pmatrix}
    g_{\mu\nu}+e^{2\phi} A_\mu A_\nu & e^{2\phi}A_\mu \\
    e^{2\phi} A_\nu  & e^{2\phi}
\end{pmatrix},
\end{equation}
which it is helpful to expand as
\begin{align}
ds^2
&= e^{-2\phi/3}
g_{\mu\nu} dx^\mu dx^\nu
+e^{4\phi /3}(dx^{10} + A_\mu dx^\mu)^2\\
&= (e^{-\phi/3}e^a)(e^{-\phi/3} e^b)\eta_{ab}+[e^{2\phi/3}(dx^{10}+A_\mu dx^\mu)]^2 .
\end{align}
Let $\hat E^A=\hat E^A{}_B dx^B$ denote the 11d frame and $e^a=e^a{}_\mu dx^\mu$ the 10d frame. Reading off,
\begin{equation}
\hat E^a = e^{-\phi/3}e^a = e^{-\phi/3}e^a{}_\mu  dx^\mu,\qquad
\hat E^{10}=e^{2\phi/3}(dx^{10}+A_\mu dx^\mu).
\end{equation}

\paragraph{A note on the hats.}
The indices $0,\dots,9$ appear both in 11d and 10d, so for a contracted object such as a scalar $T$ there is no ambiguity, but when indices are displayed there is: $T_{ab}$ may refer to an 11d or a 10d tensor, and the two differ because the frames differ by $e^{-\phi/3}$. Concretely,
\begin{equation}
T =\hat T_{ab} \hat E^a \hat E^b=T_{ab} e^a e^b
\;\implies\;
\hat T_{ab}=e^{2\phi/3}T_{ab},\qquad \hat \eta_{bc} = e^{2\phi/3}\eta_{bc}.
\end{equation}
We use a hat for 11d frame objects and the corresponding plain symbol for the 10d ones throughout; likewise $\hat\partial_a\equiv \hat E^\mu{}_a\partial_\mu=e^{\phi/3}e^\mu{}_a\partial_\mu=e^{\phi/3}\partial_a$.

\paragraph{Spin connection.}
The torsion-free Cartan equations split as
\begin{equation}
d\hat E^a+\hat W^a{}_b\wedge \hat E^b + \hat W^a{}_{10} \wedge \hat E^{10} = 0,\qquad
d\hat E^{10}+ \hat W^{10}{}_{a}\wedge \hat E^a = 0 .
\end{equation}
Solving the second equation first,
\begin{align}
d\hat E^{10}
&= d(e^{2\phi/3})\wedge (A+dx^{10}) + e^{2\phi/3}  dA\\
&=\tfrac{2}{3} d\phi\wedge \hat E^{10} + \tfrac{1}{2}e^{2\phi/3}  \hat F_{ab}  \hat E^a \wedge \hat E^b\\
&=-\left[
\tfrac{2}{3}\hat\partial_b\phi  \hat E^{10}+\tfrac{1}{2}e^{2\phi/3}\hat F_{ba}\hat E^a
\right]\wedge \hat E^b ,
\end{align}
with $\hat F_{ab}=e^{2\phi/3}F_{ab}$ and $F_{ab}=e^\mu{}_a e^\nu{}_b F_{\mu\nu}$, $F=dA$. Hence
\begin{align}
\hat W_{b,10}
&=
\underbrace{\left(-\tfrac{2}{3}e^{\phi/3}\partial_b\phi\right)}_{\hat W_{10,b,10}} \hat E^{10}
+ \underbrace{\left(\tfrac{1}{2}e^{4\phi/3}F_{ab}\right)}_{\hat W_{a,b,10}}\hat E^a .
\end{align}
For the first Cartan equation,
\begin{align}
d\hat E^a
&= -\tfrac{1}{3}d\phi \wedge \hat E^a+ e^{-\phi/3} de^a
=-\left[w^a{}_b-\tfrac{1}{3}\hat \partial_b\phi  \hat E^a\right]\wedge \hat E^b ,
\end{align}
so that, substituting $\hat W^{10}{}_b$ and fixing $\hat W^a{}_b$ by antisymmetry,
\begin{align}
\hat W_{bc}
&=\underbrace{e^{\phi/3}\left[w_{abc}+\tfrac{2}{3}e^\nu{}_{[b} \eta_{c]a}\partial_\nu \phi\right]}_{\hat W_{abc}} \hat E^a
+\underbrace{\left[-\tfrac{1}{2}e^{4\phi/3}F_{bc}\right]}_{\hat W_{10,bc}}\hat E^{10}.
\end{align}
Every component agrees with eq.~(22.19) of \cite{Ortin:2015hya}.
Contracting with gamma matrices (and using $\{\Gamma^b,\Gamma^{10}\}=0$, so $\Gamma^{b,10}=\Gamma^b\Gamma^{10}$),
\begin{align}
\hat W_{aBC}\Gamma^{BC}
&=e^{\phi/3}\left[w_{abc}\Gamma^{bc}-\tfrac{2}{3}\Gamma_a{}^{b}\partial_b\phi\right]
+e^{4\phi/3}F_{ab}\Gamma^{b}\Gamma^{10},\\
\hat W_{10,BC}\Gamma^{BC}
&=-\tfrac{1}{2}e^{4\phi/3}F_{bc}\Gamma^{bc}-\tfrac{4}{3}e^{\phi/3}\partial_\mu\phi  \Gamma^\mu \Gamma^{10},
\end{align}
matching eqs.~(8.44)--(8.45) of \cite{Becker:2006dvp}.

\paragraph{Covariant derivatives.}
The 11d covariant derivative acting on a spinor (taken independent of $x^{10}$) then gives
\begin{align}
\hat\nabla_{10} \epsilon
&=\tfrac{1}{4}\hat W_{10,BC}\Gamma^{BC}\epsilon
=\tfrac{1}{4}e^{\phi/3}\left[
-\tfrac{1}{2}e^{\phi}F_{bc}\Gamma^{bc}-\tfrac{4}{3}\partial_\mu\phi  \Gamma^\mu \Gamma^{10}
\right]\epsilon,\\
\hat\nabla_a \epsilon
&=e^{\phi/3}e^\mu{}_a
\left[
\nabla_\mu
-\tfrac{1}{6}\Gamma_\mu{}^\nu\partial_\nu\phi+\tfrac{1}{4}e^\phi F_{\mu\nu}\Gamma^\nu\Gamma^{10}
\right]\epsilon ,
\end{align}
with $\nabla_\mu=\partial_\mu+\tfrac14 w_{\mu bc}\Gamma^{bc}$ the 10d spinor covariant derivative.

\paragraph{Reduction of the Killing spinor condition.}
A Killing spinor of 11d gravity is covariantly constant, $\hat\nabla_{10}\epsilon=\hat\nabla_a\epsilon=0$. From $\hat\nabla_{10}\epsilon=0$, and using $(\Gamma^{10})^2=1$,
\begin{equation}
\Gamma^b\partial_b\phi \epsilon=\tfrac{3}{8}e^\phi F_{bc}\Gamma^{bc}\Gamma^{10}\epsilon .
\label{eq:dilatino-relation}
\end{equation}
We now eliminate $\partial\phi$ from $\hat\nabla_a\epsilon=0$. Writing $\Gamma_\mu{}^\nu\partial_\nu\phi=\Gamma_\mu(\Gamma^\nu\partial_\nu\phi)-\partial_\mu\phi$ and inserting \eqref{eq:dilatino-relation},
\begin{equation}
\left(-\tfrac{1}{6}\Gamma_\mu{}^\nu\partial_\nu\phi\right)\epsilon
=
\left(\tfrac{1}{6}\partial_\mu\phi-\tfrac{1}{8}e^\phi F_{\mu\rho}\Gamma^\rho \Gamma^{10}-\tfrac{1}{16}e^\phi F_{\nu\rho}\Gamma_\mu{}^{\nu\rho}\Gamma^{10}\right)\epsilon ,
\end{equation}
where we used $\Gamma_\mu\Gamma^{\nu\rho}=\Gamma_\mu{}^{\nu\rho}+2\delta_\mu^{[\nu}\Gamma^{\rho]}$ and $F_{\nu\rho} 2\delta_\mu^{[\nu}\Gamma^{\rho]}=2F_{\mu\nu}\Gamma^\nu$. Substituting into $\hat\nabla_a\epsilon=0$, the two $F_{\mu\nu}\Gamma^\nu\Gamma^{10}$ terms combine ($-\tfrac18+\tfrac14=+\tfrac18$) to give
\begin{equation}
\nabla_a\epsilon=e^{\phi/3}e^\mu{}_a
\left(
\nabla_\mu+\tfrac{1}{6}\partial_\mu\phi
+\tfrac{1}{8}e^\phi F_{\mu\nu}\Gamma^\nu\Gamma^{10}-\tfrac{1}{16}e^\phi F_{\nu\rho}\Gamma_\mu{}^{\nu\rho}\Gamma^{10}
\right)\epsilon=0 .
\label{eq:assembled}
\end{equation}

\paragraph{Final gamma-matrix reduction.}
The rank-one term $\tfrac18 F_{\mu\nu}\Gamma^\nu$ does not stand alone: it recombines with the rank-three term into a single flux-slash acting on $\Gamma_\mu$. Using $\Gamma^{\nu\rho}\Gamma_\mu=\Gamma_\mu{}^{\nu\rho}-2\delta_\mu^{[\nu}\Gamma^{\rho]}$ and again $F_{\nu\rho} 2\delta_\mu^{[\nu}\Gamma^{\rho]}=2F_{\mu\nu}\Gamma^\nu$,
\begin{align}
\tfrac{1}{8}F_{\mu\nu}\Gamma^\nu-\tfrac{1}{16}F_{\nu\rho}\Gamma_\mu{}^{\nu\rho}
&=-\tfrac{1}{16}F_{\nu\rho}\left(\Gamma_\mu{}^{\nu\rho}-2\delta_\mu^{[\nu}\Gamma^{\rho]}\right)\\
&=-\tfrac{1}{16}F_{\nu\rho} \Gamma^{\nu\rho}\Gamma_\mu .
\end{align}
Finally, redefining $\epsilon\to e^{-\phi/6}\epsilon$ removes the $\tfrac16\partial_\mu\phi$ term, leaving
\begin{equation}
\left(
\nabla_\mu -\tfrac{1}{16}e^\phi F_{\nu\rho} \Gamma^{\nu\rho}\Gamma_\mu  \Gamma^{10}
\right)\epsilon=0 ,
\end{equation}
which is precisely the type IIA gravitino variation \eqref{IIA Killing spinor equation}. Thus covariantly-constant spinors in 11d are Killing spinors of 10d type IIA in the gravity, KK-vector and dilaton sector.
